\documentclass[multi=tikzpicture,border=7pt]{standalone}
\usepackage{ctex,amsmath,bm}
\usepackage{tikz,pgfplots}
\pgfplotsset{compat=1.18}
\usetikzlibrary{arrows.meta,calc,angles,quotes,patterns,positioning,decorations.pathmorphing,3d}
\tikzset{wire/.style={brown!70!black,very thick},guide/.style={dashed,gray!60},force/.style={->,>=Stealth,thick,red!70!black},velocity/.style={->,>=Stealth,thick,blue!70!black},axis/.style={->,>=Stealth,thick},dim/.style={<->,>=Stealth,thin},body/.style={draw,thick,fill=blue!7},panel/.style={anchor=west,font=\bfseries},hidden/.style={dashed,gray!58,thick},orbit/.style={very thick,blue!70!black}}
\tikzset{vec/.style={->,>=Stealth,very thick},wall/.style={pattern=north east lines,pattern color=gray},nodept/.style={circle,fill=white,draw,inner sep=1.5pt}}
\begin{document}
% 1 question-01
\begin{tikzpicture}[scale=.9]
  \node at (-2.35,2.95){(a)}; \node at (3.55,2.95){(b)};
  % (a)：圆心位于自由端正下方，圆弧下端固定
  \draw[very thick] (-2.35,2.35) arc (90:0:2.35);
  \draw[very thick] (0,0)--(.48,0);
  \draw[wall] (.48,-.42) rectangle (.88,.42);
  \draw[thick] (-2.35,2.35)--(-2.35,1.25);
  \fill[gray!30] (-2.35,.88) circle(.37); \node[left=3pt] at (-2.35,.88){$m$};
  % (b)：圆心位于自由端右侧，圆弧下端固定
  \draw[very thick] (2.15,2.35) arc (180:270:2.35);
  \draw[very thick] (4.5,0)--(4.98,0);
  \draw[wall] (4.98,-.42) rectangle (5.38,.42);
  \draw[thick] (2.15,2.35)--(2.15,1.25);
  \fill[gray!30] (2.15,.88) circle(.37); \node[left=3pt] at (2.15,.88){$m$};
  \node[align=center] at (1.45,-1.15){两根同样的四分之一圆弧\\自由端承受相同竖直载荷};
\end{tikzpicture}

% 2 solution-01
\begin{tikzpicture}[x=1cm,y=1cm,font=\small]

\begin{scope}[shift={(-3,0)}]\draw[wire](0,0)arc[start angle=270,end angle=180,radius=2];\draw[guide](0,2)--(0,0) (0,2)--(-2,2);\draw[force](0,0)--(0,-1.2)node[right]{$P$};\draw[guide](0,2)--(-1.4142,.5858);\node at(-.5,1){$\varphi$};\node at(-.5,-1.9){(a) $M_a=PR\sin\varphi$};\end{scope}
\begin{scope}[shift={(3,0)}]\draw[wire](-2,0)arc[start angle=180,end angle=90,radius=2];\draw[guide](0,0)--(-2,0) (0,0)--(0,2);\draw[force](-2,0)--(-2,-1.2)node[left]{$P$};\draw[guide](0,0)--(-1.4142,1.4142);\node at(-.65,.4){$\varphi$};\node at(-.5,-1.9){(b) $M_b=PR(1-\cos\varphi)$};\end{scope}
\node at(-.8,-2.8){角度均从自由端起算；$ds=R\,d\varphi$，$0\le\varphi\le\pi/2$};
\end{tikzpicture}

% 3 question-02
\begin{tikzpicture}[x=1cm,y=1cm,font=\small]

\begin{scope}[shift={(-2.8,0)}]\draw[body,fill=none](0,0)circle(1.7);\draw[thick](0,0)circle(1.35);\node at(0,0){$\bm B:\otimes$};\draw[dim](0,-2.2)--(1.7,-2.2)node[below,midway]{$R$};\draw[dim](1.35,0)--(1.7,0);\node[right]at(1.8,.3){$d$};\node at(0,2.2){外板 $+Q$，内板 $-Q$};\end{scope}
\begin{scope}[shift={(2.5,0)}]\draw[body](-1.5,-1.2)rectangle(1.5,1.2);\draw[thick](-1.5,1.2)--(-1.1,3) (0,1.2)--(0,3) (1.5,1.2)--(1.1,3);\draw[thick](-1.3,3)--(1.3,3);\draw[dim](2,-1.2)--(2,1.2)node[right,midway]{$l$};\node at(0,-2.2){三根绝缘线悬挂；允许绕轴转动};\end{scope}
\end{tikzpicture}

% 4 solution-02
\begin{tikzpicture}[x=1cm,y=1cm,font=\small]

\foreach\x in{-3,3}{\begin{scope}[shift={(\x,0)}]\draw[body,fill=none](0,0)circle(1.7);\draw[thick](0,0)circle(1.2);\node at(0,0){$\bm B:\otimes$};\end{scope}}
\draw[velocity](-1.3,0)--(-1.8,0)node[above]{$dq>0$};\draw[force](-1.5,0)--(-1.5,-1)node[right]{$d\bm F_B$};\node at(-3,-2.3){放电：向内电流，磁力矩顺时针};
\draw[force](4.7,0)--(4.7,-1)node[right]{$+Q$};\draw[force](4.2,0)--(4.2,.9)node[left]{$-Q$};\node at(3,-2.3){撤场：两板力矩反向，外板半径较大};
\node at(0,-3.1){两过程总角冲量相同：$\Delta L_z=-QB(R^2-(R-d)^2)/2$};
\end{tikzpicture}

% 5 question-03
\begin{tikzpicture}[x=1cm,y=1cm,font=\small]

\draw[thick](-3,0)--(3,0);\draw[body](0,1)circle(1);\draw[guide](0,0)--(0,1);\draw[velocity](-2.2,2.7)--(-.7,1.7)node[left]{$\bm v_0$};\draw[velocity](.8,1.6)--(2.2,2.7)node[right]{$\bm v$};\draw[velocity](.7,1.714)arc[start angle=45,end angle=-40,radius=1];\node[right]at(1.3,1.2){$\omega_0$};\node at(0,-.7){来球向下，离地后向上；切向冲量改变平动与自转};
\end{tikzpicture}

% 6 solution-03
\begin{tikzpicture}[x=1cm,y=1cm,font=\small]

\draw[thick](-3,0)--(3,0);\draw[body](0,1)circle(1);\draw[force](0,0)--(1.8,0)node[right]{$J$（正向）};\draw[force](0,0)--(0,2.5)node[above]{$J_n$};\draw[velocity](0,1)--(1.7,1)node[right]{$v$};\draw[guide](0,1)--(0,0);\node at(0,-.8){$m(v-v_0)=J$，$I(\omega-\omega_0)=-RJ$};\node at(0,-1.5){$u=v_0-R\omega_0$ 决定 $J$ 的符号；不是始终向右};
\end{tikzpicture}

% 7 question-04
\begin{tikzpicture}[x=1cm,y=1cm,font=\small]

\draw[body](-2.5,-.3)rectangle(-1.7,.3);\draw[body](1.7,-.3)rectangle(2.5,.3);\draw[thick](-3,-.3)--(3,-.3);\draw[wire](-1.7,0)arc[start angle=180,end angle=0,radius=1.7];\node at(0,2.3){原题声称梁恰为半圆};\draw[dim](-1.7,-.8)--(1.7,-.8)node[below,midway]{$2R$};\node at(0,-1.6){若仅端点水平接触，局部弯矩条件与恒曲率矛盾};
\end{tikzpicture}

% 8 solution-04
\begin{tikzpicture}[x=1cm,y=1cm,font=\small]

\begin{scope}[shift={(-2.8,0)}]\draw[wire](-1.5,0)arc[start angle=180,end angle=0,radius=1.5];\draw[force](-1.5,0)--(-2.5,0)node[left]{$F$};\draw[force](1.5,0)--(2.5,0)node[right]{$F$};\draw[dim](-1.5,-.6)--(1.5,-.6)node[below,midway]{$s(\rho)$};\node at(0,2.1){附加圆弧形状约束};\end{scope}
\begin{scope}[shift={(3,0)}]\draw[axis](0,0)--(2.5,0)node[right]{$s$};\draw[axis](0,0)--(0,2.3)node[above]{$M$};\draw[force](0,1.4)--(2,1.4)node[above]{$YI/R$};\draw[orbit]plot[domain=0:2,samples=50](\x,{1.5*sin(90*\x)});\node at(1,-.9){端点力给 $M=Fy(s)$，端点 $M=0$};\end{scope}
\node at(0,-1.9){$F=\pi Yad^3/(24R^2)$ 是圆弧约束族的广义力，非原题唯一反力};
\end{tikzpicture}

% 9 question-05
\begin{tikzpicture}[x=1cm,y=1cm,font=\small]

\draw[thick](-2,-2)--(-2,2) (2,-2)--(2,2);\coordinate(O)at(0,3);\coordinate(P)at(.9,.5);\draw[wire](O)--(P)node[midway,right]{$L$};\draw[body](P)circle(.16);\draw[guide](0,3)--(0,.34);\draw[force](P)--++(1,0)node[right]{$qV/d$};\draw[force](P)--++(0,-1.2)node[left]{$Mg$};\draw[force](P)--++(-.45,1.25)node[left]{$T$};\draw[dim](-2,-2.5)--(2,-2.5)node[below,midway]{$d$};\draw[dim](-2.6,-2)--(-2.6,2)node[left,midway]{$h$};\node at(0,-3.3){图示平衡；未带电时球在两板中央、高度 $h/2$};
\end{tikzpicture}

% 10 solution-05
\begin{tikzpicture}[x=1cm,y=1cm,font=\small]
\draw[axis](0,0)--(7,0)node[right]{$\omega t$};\draw[axis](0,0)--(0,3.7)node[above]{$x\omega^2/(qV_0/Md)$};\draw[orbit]plot[smooth]coordinates{(0.00000,0.00000) (0.03927,0.00002) (0.07854,0.00016) (0.11781,0.00054) (0.15708,0.00129) (0.19635,0.00250) (0.23562,0.00431) (0.27489,0.00682) (0.31416,0.01013) (0.35343,0.01435) (0.39270,0.01957) (0.43197,0.02588) (0.47124,0.03337) (0.51051,0.04209) (0.54978,0.05214) (0.58905,0.06355) (0.62832,0.07640) (0.66759,0.09070) (0.70686,0.10651) (0.74613,0.12383) (0.78540,0.14270) (0.82467,0.16310) (0.86394,0.18505) (0.90321,0.20851) (0.94248,0.23347) (0.98175,0.25990) (1.02102,0.28776) (1.06029,0.31698) (1.09956,0.34752) (1.13883,0.37931) (1.17810,0.41227) (1.21737,0.44632) (1.25664,0.48137) (1.29591,0.51733) (1.33518,0.55409) (1.37445,0.59155) (1.41372,0.62960) (1.45299,0.66813) (1.49226,0.70702) (1.53153,0.74615) (1.57080,0.78540) (1.61007,0.82465) (1.64934,0.86378) (1.68861,0.90266) (1.72788,0.94119) (1.76715,0.97924) (1.80642,1.01671) (1.84569,1.05347) (1.88496,1.08942) (1.92423,1.12447) (1.96350,1.15852) (2.00277,1.19148) (2.04204,1.22327) (2.08131,1.25381) (2.12058,1.28304) (2.15984,1.31089) (2.19911,1.33732) (2.23838,1.36228) (2.27765,1.38575) (2.31692,1.40769) (2.35619,1.42810) (2.39546,1.44696) (2.43473,1.46429) (2.47400,1.48009) (2.51327,1.49440) (2.55254,1.50724) (2.59181,1.51866) (2.63108,1.52870) (2.67035,1.53743) (2.70962,1.54491) (2.74889,1.55122) (2.78816,1.55644) (2.82743,1.56066) (2.86670,1.56398) (2.90597,1.56648) (2.94524,1.56829) (2.98451,1.56951) (3.02378,1.57025) (3.06305,1.57064) (3.10232,1.57078) (3.14159,1.57080) (3.18086,1.57082) (3.22013,1.57096) (3.25940,1.57134) (3.29867,1.57208) (3.33794,1.57330) (3.37721,1.57511) (3.41648,1.57762) (3.45575,1.58093) (3.49502,1.58515) (3.53429,1.59037) (3.57356,1.59668) (3.61283,1.60416) (3.65210,1.61289) (3.69137,1.62293) (3.73064,1.63435) (3.76991,1.64719) (3.80918,1.66150) (3.84845,1.67730) (3.88772,1.69463) (3.92699,1.71350) (3.96626,1.73390) (4.00553,1.75584) (4.04480,1.77931) (4.08407,1.80427) (4.12334,1.83070) (4.16261,1.85855) (4.20188,1.88778) (4.24115,1.91832) (4.28042,1.95011) (4.31969,1.98307) (4.35896,2.01712) (4.39823,2.05217) (4.43750,2.08813) (4.47677,2.12489) (4.51604,2.16235) (4.55531,2.20040) (4.59458,2.23893) (4.63385,2.27782) (4.67312,2.31694) (4.71239,2.35619) (4.75166,2.39544) (4.79093,2.43457) (4.83020,2.47346) (4.86947,2.51199) (4.90874,2.55004) (4.94801,2.58750) (4.98728,2.62426) (5.02655,2.66022) (5.06582,2.69527) (5.10509,2.72932) (5.14436,2.76228) (5.18363,2.79407) (5.22290,2.82461) (5.26217,2.85384) (5.30144,2.88169) (5.34071,2.90812) (5.37998,2.93308) (5.41925,2.95655) (5.45852,2.97849) (5.49779,2.99889) (5.53706,3.01776) (5.57633,3.03509) (5.61560,3.05089) (5.65487,3.06520) (5.69414,3.07804) (5.73341,3.08945) (5.77268,3.09950) (5.81195,3.10823) (5.85122,3.11571) (5.89049,3.12202) (5.92976,3.12724) (5.96903,3.13146) (6.00830,3.13477) (6.04757,3.13728) (6.08684,3.13909) (6.12611,3.14031) (6.16538,3.14105) (6.20465,3.14143) (6.24392,3.14157) (6.28319,3.14159)};\node[below]at(3.14159,0){$2\pi$};\node[below]at(6.28318,0){$4\pi$};\draw[guide](3.14159,0)--(3.14159,1.5708);\node at(3.2,-.9){$x\propto\omega t-\sin\omega t$，始终不减；驻点不代表返回};
\end{tikzpicture}

% 11 question-06
\begin{tikzpicture}[x=1cm,y=1cm,font=\small]
\draw[wire](0.00000,3.46410)--(-0.07407,3.33580);
\draw[wire](-0.07407,3.33580)--(0.07407,3.33580);
\draw[wire](0.07407,3.33580)--(0.00000,3.46410);
\draw[wire](-0.07407,3.33580)--(-0.14815,3.20750);
\draw[wire](-0.22222,3.07920)--(-0.07407,3.07920);
\draw[wire](-0.07407,3.07920)--(-0.14815,3.20750);
\draw[wire](-0.14815,3.20750)--(-0.22222,3.07920);
\draw[wire](-0.07407,3.07920)--(0.07407,3.07920);
\draw[wire](0.22222,3.07920)--(0.14815,3.20750);
\draw[wire](0.14815,3.20750)--(0.07407,3.07920);
\draw[wire](0.07407,3.07920)--(0.22222,3.07920);
\draw[wire](0.14815,3.20750)--(0.07407,3.33580);
\draw[wire](-0.22222,3.07920)--(-0.44444,2.69430);
\draw[wire](-0.66667,2.30940)--(-0.51852,2.30940);
\draw[wire](-0.51852,2.30940)--(-0.59259,2.43770);
\draw[wire](-0.59259,2.43770)--(-0.66667,2.30940);
\draw[wire](-0.51852,2.30940)--(-0.37037,2.30940);
\draw[wire](-0.22222,2.30940)--(-0.29630,2.43770);
\draw[wire](-0.29630,2.43770)--(-0.37037,2.30940);
\draw[wire](-0.37037,2.30940)--(-0.22222,2.30940);
\draw[wire](-0.29630,2.43770)--(-0.37037,2.56600);
\draw[wire](-0.44444,2.69430)--(-0.51852,2.56600);
\draw[wire](-0.51852,2.56600)--(-0.37037,2.56600);
\draw[wire](-0.37037,2.56600)--(-0.44444,2.69430);
\draw[wire](-0.51852,2.56600)--(-0.59259,2.43770);
\draw[wire](-0.22222,2.30940)--(0.22222,2.30940);
\draw[wire](0.66667,2.30940)--(0.59259,2.43770);
\draw[wire](0.59259,2.43770)--(0.51852,2.30940);
\draw[wire](0.51852,2.30940)--(0.66667,2.30940);
\draw[wire](0.59259,2.43770)--(0.51852,2.56600);
\draw[wire](0.44444,2.69430)--(0.37037,2.56600);
\draw[wire](0.37037,2.56600)--(0.51852,2.56600);
\draw[wire](0.51852,2.56600)--(0.44444,2.69430);
\draw[wire](0.37037,2.56600)--(0.29630,2.43770);
\draw[wire](0.22222,2.30940)--(0.37037,2.30940);
\draw[wire](0.37037,2.30940)--(0.29630,2.43770);
\draw[wire](0.29630,2.43770)--(0.22222,2.30940);
\draw[wire](0.37037,2.30940)--(0.51852,2.30940);
\draw[wire](0.44444,2.69430)--(0.22222,3.07920);
\draw[wire](-0.66667,2.30940)--(-1.33333,1.15470);
\draw[wire](-2.00000,0.00000)--(-1.85185,0.00000);
\draw[wire](-1.85185,0.00000)--(-1.92593,0.12830);
\draw[wire](-1.92593,0.12830)--(-2.00000,0.00000);
\draw[wire](-1.85185,0.00000)--(-1.70370,0.00000);
\draw[wire](-1.55556,0.00000)--(-1.62963,0.12830);
\draw[wire](-1.62963,0.12830)--(-1.70370,0.00000);
\draw[wire](-1.70370,0.00000)--(-1.55556,0.00000);
\draw[wire](-1.62963,0.12830)--(-1.70370,0.25660);
\draw[wire](-1.77778,0.38490)--(-1.85185,0.25660);
\draw[wire](-1.85185,0.25660)--(-1.70370,0.25660);
\draw[wire](-1.70370,0.25660)--(-1.77778,0.38490);
\draw[wire](-1.85185,0.25660)--(-1.92593,0.12830);
\draw[wire](-1.55556,0.00000)--(-1.11111,0.00000);
\draw[wire](-0.66667,0.00000)--(-0.74074,0.12830);
\draw[wire](-0.74074,0.12830)--(-0.81481,0.00000);
\draw[wire](-0.81481,0.00000)--(-0.66667,0.00000);
\draw[wire](-0.74074,0.12830)--(-0.81481,0.25660);
\draw[wire](-0.88889,0.38490)--(-0.96296,0.25660);
\draw[wire](-0.96296,0.25660)--(-0.81481,0.25660);
\draw[wire](-0.81481,0.25660)--(-0.88889,0.38490);
\draw[wire](-0.96296,0.25660)--(-1.03704,0.12830);
\draw[wire](-1.11111,0.00000)--(-0.96296,0.00000);
\draw[wire](-0.96296,0.00000)--(-1.03704,0.12830);
\draw[wire](-1.03704,0.12830)--(-1.11111,0.00000);
\draw[wire](-0.96296,0.00000)--(-0.81481,0.00000);
\draw[wire](-0.88889,0.38490)--(-1.11111,0.76980);
\draw[wire](-1.33333,1.15470)--(-1.40741,1.02640);
\draw[wire](-1.40741,1.02640)--(-1.25926,1.02640);
\draw[wire](-1.25926,1.02640)--(-1.33333,1.15470);
\draw[wire](-1.40741,1.02640)--(-1.48148,0.89810);
\draw[wire](-1.55556,0.76980)--(-1.40741,0.76980);
\draw[wire](-1.40741,0.76980)--(-1.48148,0.89810);
\draw[wire](-1.48148,0.89810)--(-1.55556,0.76980);
\draw[wire](-1.40741,0.76980)--(-1.25926,0.76980);
\draw[wire](-1.11111,0.76980)--(-1.18519,0.89810);
\draw[wire](-1.18519,0.89810)--(-1.25926,0.76980);
\draw[wire](-1.25926,0.76980)--(-1.11111,0.76980);
\draw[wire](-1.18519,0.89810)--(-1.25926,1.02640);
\draw[wire](-1.55556,0.76980)--(-1.77778,0.38490);
\draw[wire](-0.66667,0.00000)--(0.66667,0.00000);
\draw[wire](2.00000,0.00000)--(1.92593,0.12830);
\draw[wire](1.92593,0.12830)--(1.85185,0.00000);
\draw[wire](1.85185,0.00000)--(2.00000,0.00000);
\draw[wire](1.92593,0.12830)--(1.85185,0.25660);
\draw[wire](1.77778,0.38490)--(1.70370,0.25660);
\draw[wire](1.70370,0.25660)--(1.85185,0.25660);
\draw[wire](1.85185,0.25660)--(1.77778,0.38490);
\draw[wire](1.70370,0.25660)--(1.62963,0.12830);
\draw[wire](1.55556,0.00000)--(1.70370,0.00000);
\draw[wire](1.70370,0.00000)--(1.62963,0.12830);
\draw[wire](1.62963,0.12830)--(1.55556,0.00000);
\draw[wire](1.70370,0.00000)--(1.85185,0.00000);
\draw[wire](1.77778,0.38490)--(1.55556,0.76980);
\draw[wire](1.33333,1.15470)--(1.25926,1.02640);
\draw[wire](1.25926,1.02640)--(1.40741,1.02640);
\draw[wire](1.40741,1.02640)--(1.33333,1.15470);
\draw[wire](1.25926,1.02640)--(1.18519,0.89810);
\draw[wire](1.11111,0.76980)--(1.25926,0.76980);
\draw[wire](1.25926,0.76980)--(1.18519,0.89810);
\draw[wire](1.18519,0.89810)--(1.11111,0.76980);
\draw[wire](1.25926,0.76980)--(1.40741,0.76980);
\draw[wire](1.55556,0.76980)--(1.48148,0.89810);
\draw[wire](1.48148,0.89810)--(1.40741,0.76980);
\draw[wire](1.40741,0.76980)--(1.55556,0.76980);
\draw[wire](1.48148,0.89810)--(1.40741,1.02640);
\draw[wire](1.11111,0.76980)--(0.88889,0.38490);
\draw[wire](0.66667,0.00000)--(0.81481,0.00000);
\draw[wire](0.81481,0.00000)--(0.74074,0.12830);
\draw[wire](0.74074,0.12830)--(0.66667,0.00000);
\draw[wire](0.81481,0.00000)--(0.96296,0.00000);
\draw[wire](1.11111,0.00000)--(1.03704,0.12830);
\draw[wire](1.03704,0.12830)--(0.96296,0.00000);
\draw[wire](0.96296,0.00000)--(1.11111,0.00000);
\draw[wire](1.03704,0.12830)--(0.96296,0.25660);
\draw[wire](0.88889,0.38490)--(0.81481,0.25660);
\draw[wire](0.81481,0.25660)--(0.96296,0.25660);
\draw[wire](0.96296,0.25660)--(0.88889,0.38490);
\draw[wire](0.81481,0.25660)--(0.74074,0.12830);
\draw[wire](1.11111,0.00000)--(1.55556,0.00000);
\draw[wire](1.33333,1.15470)--(0.66667,2.30940);
\fill(0.00000,3.46410)circle(1.8pt)node[above]{$ A $};
\fill(-2.00000,0.00000)circle(1.8pt)node[above]{$ B $};
\fill(2.00000,0.00000)circle(1.8pt)node[above]{$ C $};
\draw[velocity](-.18,3.15)--(-1.8,.32);\node[left]at(-1.5,1.7){外边 $A\to B$};\node at(2.8,2){$\dot{\bm B}:\otimes$};\node at(0,-.8){图示前三次替换；每个角落继续同样递归};
\end{tikzpicture}

% 12 solution-06
\begin{tikzpicture}[x=1cm,y=1cm,font=\small]
\coordinate(N0)at(0.00000,2.00000);\fill(N0)circle(1.5pt);
\coordinate(N1)at(-1.73205,1.00000);\fill(N1)circle(1.5pt);
\coordinate(N2)at(-1.73205,-1.00000);\fill(N2)circle(1.5pt);
\coordinate(N3)at(-0.00000,-2.00000);\fill(N3)circle(1.5pt);
\coordinate(N4)at(1.73205,-1.00000);\fill(N4)circle(1.5pt);
\coordinate(N5)at(1.73205,1.00000);\fill(N5)circle(1.5pt);
\draw[wire](N0)--(N1)node[midway,fill=white]{$R_0/6$};
\draw[wire](N1)--(N2)node[midway,fill=white]{$R_0/3$};
\draw[wire](N2)--(N3)node[midway,fill=white]{$R_0/6$};
\draw[wire](N3)--(N4)node[midway,fill=white]{$R_0/3$};
\draw[wire](N4)--(N5)node[midway,fill=white]{$R_0/6$};
\draw[wire](N5)--(N0)node[midway,fill=white]{$R_0/3$};
\draw[velocity](-2.35,1)--(-2.35,-.3)node[left]{$I_t$};\node at(0,0){$\mathcal E_{\rm eff}=7\mathcal E_0/9$};\node at(0,-2.8){$R_{\rm ring}=3R_0/2$，$I_t=14\mathcal E_0/(27R_0)$};\node at(0,-3.5){$P=P/9+(3/2)R_0I_t^2\Rightarrow P=49\mathcal E_0^2/(108R_0)$};
\end{tikzpicture}

% 13 question-07
\begin{tikzpicture}[x=1cm,y=1cm,font=\small]

\begin{scope}[shift={(-3,0)}]\draw[body](-1.2,-1)rectangle(1.2,1);\node at(0,.4){房间};\node at(0,-.2){$T,\ p$};\draw[guide,->](1.2,0)--(2,0);\node at(0,-1.6){微缝使压强与外界相等};\node at(0,1.6){(1) 定压加热};\end{scope}
\begin{scope}[shift={(3.3,0)}]\draw[body](0,1.8)circle(.5);\node at(0,1.8){$T_a$};\draw[body](-.5,-.4)rectangle(.5,.4);\node at(0,0){热机};\draw[body](0,-1.8)circle(.5);\node at(0,-1.8){$T_b$};\draw[force](0,1.3)--(0,.4)node[right,midway]{$dQ_a$};\draw[velocity](0,-.4)--(0,-1.3)node[right,midway]{$dQ_b$};\draw[force](.5,0)--(1.6,0)node[right]{$dW$};\node at(0,2.75){(2) 有限冷热物体};\node at(0,-2.75){热容 $C_a,C_b$；可逆取功至两温相等};\end{scope}
\end{tikzpicture}

% 14 solution-07
\begin{tikzpicture}[x=1cm,y=1cm,font=\small]
\draw[axis](0,0)--(3.8,0)node[right]{$T_a$};\draw[axis](0,0)--(0,3.8)node[above]{$T_b$};\draw[guide](0,0)--(3.5,3.5);\draw[orbit]plot[smooth]coordinates{(3.14286,0.28571) (3.12584,0.29849) (3.10883,0.31131) (3.09182,0.32416) (3.07481,0.33706) (3.05779,0.34999) (3.04078,0.36296) (3.02377,0.37597) (3.00676,0.38902) (2.98974,0.40211) (2.97273,0.41524) (2.95572,0.42841) (2.93871,0.44161) (2.92169,0.45486) (2.90468,0.46815) (2.88767,0.48148) (2.87065,0.49485) (2.85364,0.50826) (2.83663,0.52171) (2.81962,0.53520) (2.80260,0.54873) (2.78559,0.56231) (2.76858,0.57593) (2.75157,0.58959) (2.73455,0.60329) (2.71754,0.61703) (2.70053,0.63082) (2.68352,0.64465) (2.66650,0.65852) (2.64949,0.67243) (2.63248,0.68639) (2.61547,0.70039) (2.59845,0.71444) (2.58144,0.72853) (2.56443,0.74266) (2.54741,0.75684) (2.53040,0.77107) (2.51339,0.78533) (2.49638,0.79965) (2.47936,0.81400) (2.46235,0.82841) (2.44534,0.84286) (2.42833,0.85735) (2.41131,0.87189) (2.39430,0.88648) (2.37729,0.90112) (2.36028,0.91580) (2.34326,0.93052) (2.32625,0.94530) (2.30924,0.96012) (2.29222,0.97499) (2.27521,0.98991) (2.25820,1.00488) (2.24119,1.01989) (2.22417,1.03495) (2.20716,1.05007) (2.19015,1.06523) (2.17314,1.08044) (2.15612,1.09570) (2.13911,1.11101) (2.12210,1.12637) (2.10509,1.14178) (2.08807,1.15724) (2.07106,1.17275) (2.05405,1.18831) (2.03704,1.20393) (2.02002,1.21959) (2.00301,1.23531) (1.98600,1.25108) (1.96898,1.26690) (1.95197,1.28277) (1.93496,1.29870) (1.91795,1.31468) (1.90093,1.33072) (1.88392,1.34680) (1.86691,1.36294) (1.84990,1.37914) (1.83288,1.39539) (1.81587,1.41169) (1.79886,1.42805) (1.78185,1.44447) (1.76483,1.46094) (1.74782,1.47746) (1.73081,1.49405) (1.71379,1.51069) (1.69678,1.52738) (1.67977,1.54413) (1.66276,1.56094) (1.64574,1.57781) (1.62873,1.59474) (1.61172,1.61172)};\fill(3.14286,.28571)circle(1.8pt)node[right]{初态 $(400,300)$};\fill(1.61172,1.61172)circle(1.8pt)node[above left]{$T_f=\sqrt{T_aT_b}$};\node at(1.6,-.9){等热容示例：可逆轨迹 $T_aT_b=\mathrm{const.}$；终态两温相等};
\end{tikzpicture}

% 15 question-08
\begin{tikzpicture}
  \filldraw[fill=blue!8,very thick] (-3,0) circle(.65) node{$m$}; \filldraw[fill=blue!8,very thick] (3,0) circle(.65) node{$m$};
  \draw[very thick,decorate,decoration={coil,aspect=.35,segment length=5pt,amplitude=5pt}] (-2.35,0)--(2.35,0);
  \draw[dim] (-3,1.1)--(3,1.1) node[midway,above]{$R$}; \node at (0,-1.1){弹簧劲度系数 $k$};
  \draw[vec,red] (-4.4,-1.8)--(4.4,-1.8) node[right]{$\bm E$};
\end{tikzpicture}

% 16 solution-08
\begin{tikzpicture}
  \filldraw[fill=blue!8,very thick] (-3,0) circle(.55); \filldraw[fill=blue!8,very thick] (3,0) circle(.55);
  \draw[vec,red] (-3,0)--(-1.5,0) node[above]{$\bm p_1$}; \draw[vec,red] (3,0)--(4.5,0) node[above]{$\bm p_2$};
  \draw[vec,blue] (-2.4,-.7)--(-1,-.7) node[right]{吸引}; \draw[vec,blue] (2.4,-.7)--(1,-.7);
  \node at (0,1.5){平行外场：同轴偶极相吸};
  \filldraw[fill=blue!8,very thick] (-3,-3) circle(.55); \filldraw[fill=blue!8,very thick] (3,-3) circle(.55);
  \draw[vec,red] (-3,-3)--(-3,-1.7) node[above]{$\bm p_1$}; \draw[vec,red] (3,-3)--(3,-1.7) node[above]{$\bm p_2$};
  \draw[vec,blue] (-2.4,-3.7)--(-4,-3.7) node[left]{排斥}; \draw[vec,blue] (2.4,-3.7)--(4,-3.7);
  \node at (0,-4.65){垂直外场：并列偶极相斥，作用强度为同轴情形的一半};
\end{tikzpicture}

\end{document}
